\documentclass[main.tex]{subfiles}
\begin{document}

\section{Exercise 5 – Input capacitance of the amplifier (Figure~1.1)}

\subsection*{Given}
\[
C_1=0.1~\mathrm{pF}, \quad 
C_2=0.05~\mathrm{pF}, \quad
I_D=250~\mu\mathrm{A}, \quad
V_{\text{eff1}}=0.859~\mathrm{V}, \quad
\lambda_2=0.1~\mathrm{V^{-1}}.
\]
\[
g_{m1}=\frac{2I_D}{V_{\text{eff1}}}
=\frac{2(250~\mu\mathrm{A})}{0.859}
=0.582~\mathrm{mS}.
\]
For the PMOS load,
\[
r_{ds2}=\frac{1}{\lambda_2 I_D}
=\frac{1}{0.1\cdot250~\mu\mathrm{A}}
=40~\mathrm{k}\Omega
\quad
\big(\text{or } r_{ds2}=\tfrac{1+\lambda_2 V_{SD}}{\lambda_2 I_D}\approx 46~\mathrm{k}\Omega \text{ with } V_{SD}=1.5~\mathrm{V}\big).
\]

\subsection*{Small-signal gain without capacitors}
We define \(A_{v0}\) as the low-frequency small-signal gain obtained from the transistor-only model (capacitors open-circuit).  
At the output node, the PMOS load contributes \(r_{ds2}\) to AC ground with \(v_{gs1}=v_{in}\):
\[
\frac{v_o}{r_{ds2}} + g_{m1}v_{in} = 0
\;\Rightarrow\;
A_{v0}\equiv\frac{v_o}{v_{in}} \approx -g_{m1}r_{ds2}.
\]
so,
\[
A_{v0}=-0.582~\mathrm{mS}\times40~\mathrm{k}\Omega=-23.3.
\]

\subsection*{Miller input capacitance}
\[
C_\text{in}=C_1 + C_2\,(1+|A_{v0}|).
\]
So
\[
C_\text{in}=0.1~\mathrm{pF}+0.05~\mathrm{pF}\,(1+23.3)=1.31~\mathrm{pF}.
\]

\subsection*{Result}
\[
\boxed{C_\text{in}\approx 1.3~\mathrm{pF}}
\]
Using \(r_{ds2}\approx 46~\mathrm{k}\Omega\) instead gives \(|A_{v0}|\approx 26.7\) and \(C_\text{in}\approx 1.44~\mathrm{pF}\).

\end{document}
